\usepackage{fontspec}

% For nice Kerning
\usepackage[final]{microtype}
\usepackage{fontspec}

\setmainfont[Ligatures=TeX]{IBM Plex Sans}

\usepackage[letterpaper, margin=1in]{geometry}

\usepackage{tikz}
\usepackage{siunitx}
\usepackage{graphicx}
\usepackage{setspace}
\usepackage{anyfontsize}
\usepackage{xcolor}
\usepackage{fancyhdr}
\usepackage{verbatim}
\usepackage{enumitem}
\usepackage{xstring}
\usepackage{colortbl}
\usepackage{adjustbox}
\usepackage{caption}
\usepackage{multicol}
\usepackage{multirow}

% Counter for final counts
\usepackage{totcount}
\usepackage{ifthen}

%%% Configuring code block stuff
\usepackage{listings}

\usepackage[final]{pdfpages}

%Caption Label Color
\DeclareCaptionLabelFormat{colored}{\color{\currentFindingColor}#1 #2}
\captionsetup[lstlisting]{labelformat=colored}

%Caption Color
\DeclareCaptionFont{colored}{\color{\currentFindingColor}}
\captionsetup[lstlisting]{font={colored,small}}

\definecolor{lstgreen}{rgb}{0,0.6,0}
\definecolor{lstgray}{rgb}{0.5,0.5,0.5}
\lstset{
  basicstyle=\ttfamily\small,
  commentstyle=\color{lstgray},
  keywordstyle=\color{lstgreen},
  stringstyle=\color{blue},
  frame=l,
  framexleftmargin=3pt,
  rulecolor=\color{\currentFindingColor},
  columns=flexible,
  breakatwhitespace=false,
  keepspaces=true,
  breaklines=true,
  tabsize=2,
  captionpos=b
}
%%%%

\usepackage{tabularx}
\usepackage{xltabular}
\setlength{\extrarowheight}{5pt}
\newcolumntype{s}{>{\centering\arraybackslash}X}

\usepackage{hyperref}
\hypersetup{
    colorlinks,
    linkcolor={black},
    citecolor={blue!50!black},
    urlcolor={blue!80!black}
}

%SECTION SPACING
\usepackage{titlesec}
\titlespacing\section{0pt}{12pt plus 4pt minus 2pt}{7pt plus 2pt minus 2pt}
\titlespacing\subsection{0pt}{12pt plus 4pt minus 2pt}{7pt plus 2pt minus 2pt}
\titlespacing\subsubsection{0pt}{12pt plus 4pt minus 2pt}{7pt plus 2pt minus 2pt}

%REMOVE SECTION NUMBERING BUT KEEP IN TOC
% Commented out based on feedback
%\titleformat{\section}{\normalfont\LARGE\bfseries}{}{0pt}{}
%\titleformat{\subsection}{\normalfont\large\bfseries}{}{0pt}{}
\titleformat{\subsubsection}{\color{gray}\normalfont\large}{}{0pt}{}

%ITEMIZE STUFF
\renewcommand\labelitemi{{$\bullet$}}                   %Make itemize bullet points bigger
\setitemize{itemsep=-10pt}                              %Make itemize spacing lesser
\setenumerate{itemsep=-10pt}                            %Make enumerate spacing lesser

%One more color for that one table
\definecolor{darkBlue}{HTML}{002060}

%COLORS
\definecolor{infoBlue}{HTML}{007CBF}
\definecolor{lowGreen}{HTML}{27AE4D}
\definecolor{mediumOrange}{HTML}{f5890f}
\definecolor{highRed}{HTML}{FF3030}
\definecolor{criticalPurple}{HTML}{600080}

%COLORS for background
\definecolor{infoBlueBackground}{HTML}{1995D8}
\definecolor{lowGreenBackground}{HTML}{40C766}
\definecolor{mediumOrangeBackground}{HTML}{FFA228}
\definecolor{highRedBackground}{HTML}{FF3030}
\definecolor{criticalPurpleBackground}{HTML}{600080}

%Current Finding Color, changes to color coordinate findings
\newcommand{\currentFindingColor}{black}
\newcommand{\currentFindingBackground}{black}
%Some other variables related to figures and evidence
\newcommand{\imageWidth}{6in}
\newcommand{\imageBorderWidth}{1pt}

%Gap to close the table
\newcommand{\closeTableGap}[0]{\vspace{-0.45cm}}
\newcommand{\postTableGap}[0]{\vspace{0.3cm}}


\renewcommand{\headrulewidth}{0pt}                      %Remove weird header line
\renewcommand{\contentsname}{Table of Contents}         %Change TOC Title

%CUSTOM COMMANDS
\newcommand{\leaveColor}[1]{\leavevmode\color{#1}}      %Color with leaveevmode
\newif\ifaffectedSystemsSep                             %For making affected systems table (with recursion)
\newcommand*{\affectedSystems}[1]%
    {
    \arrayrulecolor{\currentFindingColor}
    \par\begin{tabularx}{\textwidth}{|s|s|s|s|}
        \multicolumn{4}{|>{\centering\arraybackslash}p{\dimexpr0.97\linewidth-2\arrayrulewidth}|}{\textbf{Affected Systems}}\\
        \hline
        \textbf{IP Address} & \textbf{Port} & \textbf{Service} & \textbf{Version}\\
        \hline
        \affectedSystemsSepfalse
        \affectedSystemsScan#1\relax\relax\relax\relax
    \end{tabularx}\par
    \arrayrulecolor{black}
    \postTableGap
}
\newcommand{\affectedSystemsScan}[4]%
  {\ifx\relax#1\empty
  \else
    \ifaffectedSystemsSep
    \else
      \affectedSystemsSeptrue
    \fi
    #1 & #2 & #3 & #4\\\hline
    \expandafter\affectedSystemsScan
  \fi
}
\newcommand\cvss[4]%                                    %For making the CVSS table
{
    \arrayrulecolor{\currentFindingColor}
    {\begin{tabularx}{\textwidth}{|s|s|s|s|s|}
        \hline
        \vspace{-13pt}\leaveColor{%
            \ifdim#1pt>8.99pt
                criticalPurple%
            \else
                \ifdim#1pt>6.99pt
                    highRed%
                \else
                    \ifdim#1pt>3.99pt
                        mediumOrange%
                    \else
                        \ifdim#1pt>0.99pt
                            lowGreen%
                        \else
                            infoBlue%
                        \fi
                    \fi
                \fi
            \fi
        } \Huge  \multirow{3}*{\textbf{#1}}
            & \multicolumn{4}{c|}{\textbf{Adjusted CVSS v3.1 Score}} \\
        \cline{2-5}
        & \textbf{Vector} & \multicolumn{3}{c|}{#4}\\
        \cline{2-5}
        & \textbf{Likelihood} & \textbf{\IfEqCase{#2}{%
            {0}{\leaveColor{infoBlue} Insignificant}%
            {1}{\leaveColor{lowGreen} Low}%
            {2}{\leaveColor{mediumOrange} Moderate}%
            {3}{\leaveColor{highRed} High}%
            {4}{\leaveColor{criticalPurple} Very High}%
        }}& \textbf{Impact} & \leaveColor{criticalPurple} \textbf{\IfEqCase{#3}{%
            {0}{\leaveColor{infoBlue} Insignificant}%
            {1}{\leaveColor{lowGreen} Tolerable}%
            {2}{\leaveColor{mediumOrange} Moderate}%
            {3}{\leaveColor{highRed} Serious}%
            {4}{\leaveColor{criticalPurple} Catastrophic}%
        }}\\
        \hline
    \end{tabularx}%
    \closeTableGap
    }
}
\newcommand\includeEvidence[1]%                         %For including evidence. centered & set width
    {\begin{center}
        \begin{adjustbox}{cframe={\currentFindingColor} \imageBorderWidth}
            \includegraphics[width=\imageWidth]{#1}  
        \end{adjustbox}
    \end{center}  
}
\newcommand\includeFigure[3]%                           %For including a figure. centered & set width
    {\begin{figure}[h!]
        \centering
        \begin{adjustbox}{cframe={\currentFindingColor} \imageBorderWidth}
            \includegraphics[width=\imageWidth]{#1}  
        \end{adjustbox}
        \captionsetup{labelfont={color=\currentFindingColor},textfont={color=\currentFindingColor}}
        \caption{#2}
        \label{#3}
    \end{figure}  
}

%REMOVE TOC NUMBERING
\usepackage{tocloft}
% Commented out based on feedback
%\renewcommand\numberline[1]{}

% RISK SUBSECTION COMMAND
\newcommand{\riskSubsection}[1]
{
    {
    \renewcommand{\arraystretch}{1.85}
    \addcontentsline{toc}{subsection}{#1}
    \begin{tabularx}{\textwidth}{>{\hsize=0.0001\hsize}X >{\hsize=0.9999\hsize}X}
    \cellcolor{\currentFindingBackground} & \raisebox{0.15em}{\huge \leaveColor{\currentFindingColor}#1}
    \end{tabularx}
    }
}
% FINDING SUBSUBSECTION COMMAND
\newcommand{\findingSubsubsection}[1]
{
    {
    \renewcommand{\arraystretch}{1.25}
    \addcontentsline{toc}{subsubsection}{\thefinding #1}
    \begin{tabularx}{\textwidth}{>{\hsize=0.0001\hsize}X >{\hsize=0.9999\hsize}X}
    \cellcolor{\currentFindingBackground} & \raisebox{0.15em}{\large \leaveColor{\currentFindingColor}\thefinding #1}
    \end{tabularx}
    \medskip
    }
}
% FINDINGS SECTION LIKE DETAILS, CONFIRMATION, ETC
\newcommand{\findingSect}[1]{
    {\color{\currentFindingColor} \large \textbf{#1}}
}


% REMEDIATION ENTRY COMMAND
%\newcommand{\remediationEntry}[1]
%{
%    {
%    \renewcommand{\arraystretch}{1.25}
%    \begin{tabularx}{\textwidth}{>{\hsize=0.0001\hsize}X >{\hsize=0.9999\hsize}X}
%    \cellcolor{\currentFindingBackground} & \raisebox{0.15em}{\large %\leaveColor{\currentFindingColor}#1}
%    \end{tabularx}
%    \medskip
%    }
%}
%\newcommand{\remediationHeader}
%{
%    {
%    %\renewcommand{\arraystretch}{1.25}
%    \begin{tabularx}{\textwidth}{| >{\hsize=0.666\hsize}X !{\vrule width 1.2pt} >{\hsize=0.333\hsize\centering\arraybackslash}X |}
%        \hline
%        \cellcolor{darkBlue}\leaveColor{white}\textbf{Vulnerability} & \cellcolor{darkBlue}\leaveColor{white}\textbf{Remediation Status}\\
%        \hline
%    \end{tabularx}
%    }
%    %\vspace{-0.4pt}
%}
%\newcommand{\remediationEntry}[2]
%{
%    \vspace{-0.4pt}
%    {
%    %\noindent
%    %\vspace{-\baselineskip}
%    %\vspace{-17.4pt}
%    \renewcommand{\arraystretch}{1.25}
%    \begin{tabularx}{\textwidth}{| >{\hsize=0.666\hsize}X | >{\hsize=0.333\hsize\centering\arraybackslash}X |}
%        \leaveColor{\currentFindingColor}\textbf{#1} & \cellcolor{\currentFindingBackground}\leaveColor{white}\textbf{#2}\\
%        \hline
%    \end{tabularx}
%    }
%}
%\newcommand{\remediationDescription}[1]
%{
%    {
%    %\noindent
%    %\vspace{-\baselineskip}
%    %\vspace{-18.4pt}
%    \renewcommand{\arraystretch}{1.25}
%    \begin{tabularx}{\textwidth}{| X |}
%        #1\\
%        \hline
%    \end{tabularx}
%    }\\
%    %\vspace{5pt}
%}

\newcommand{\remediationHeader}[1]
{
    {
    %\renewcommand{\arraystretch}{1.25}
    \begin{tabularx}{\textwidth}{| X |}
        \hline
        \cellcolor{\currentFindingBackground}\leaveColor{white}\textbf{#1}\\
        \hline
    \end{tabularx}
    }
    \vspace{-2.4pt}
}
\newcommand{\remediationEntry}[2]
{
    %\vspace{-0.4pt}
    {
    %\noindent
    \vspace{-\baselineskip}
    %\vspace{-17.4pt}
    \renewcommand{\arraystretch}{1.25}
    \begin{tabularx}{\textwidth}{| X |}
        \leaveColor{\currentFindingColor}\textbf{#1} \\
        \hline
    \end{tabularx}
    }
}
%\newcommand{\remediationDescription}[1]
%{
%    {
%    %\noindent
%    \vspace{-\baselineskip}
%    \vspace{1pt}
%    %\vspace{-18.4pt}
%    \renewcommand{\arraystretch}{1.25}
%    \begin{tabularx}{\textwidth}{| X |}
%        #1\\
%        \hline
%    \end{tabularx}
%    }\\
%    %\vspace{5pt}
%}

%For auto noindent (after tocloft)
\usepackage{parskip}        
\setlength{\parskip}{0.75\baselineskip plus 2pt}

\onehalfspacing


% For tools appendix
\newcommand{\toolref}[4]{
    \subsubsection*{#1 \hfill \small #2}
    \label{sec:#1}

    #3

    %\texttt{Source: \href{#4}{#4}}
    \texttt{\href{#4}{#4}}
}

% Reference URL and show the URL
\newcommand{\aref}[1]{
    \href{#1}{\url{#1}}
}
% Eat the period in bib
\newcommand\EatDot[1]{}
% Refernece URL in bib to do all the magic
\newcommand\bibURL[1]{
    \\ \aref{#1}\EatDot
}

\newcommand{\closeFinding}{
    \begin{center}
    \textbf{\Large END OF FINDING BLOCK}
    \end{center}
    \newpage
}

%%% Coutner for the findings

%%totcounter
\newtotcounter{totalFindings}
\newtotcounter{criticalFindings}
\newtotcounter{highFindings}
\newtotcounter{mediumFindings}
\newtotcounter{lowFindings}
\newtotcounter{informationalFindings}
%%

\newcounter{finding}
\newcounter{findingSection}
\newcommand{\findingSection}{N}
\setcounter{finding}{1}
\renewcommand{\thefinding}{\findingSection\arabic{finding}: }

\newcommand{\setFindingSection}[1]{%
    \renewcommand{\findingSection}{#1}%
    \setcounter{finding}{1}%
}
\newcommand{\incrementFinding}{%
    \stepcounter{finding}%
    \stepcounter{totalFindings}%
    \ifthenelse{\equal{\currentFindingColor}{criticalPurple}}{
        % Code for critical purple
        \stepcounter{criticalFindings}
    }{
        \ifthenelse{\equal{\currentFindingColor}{highRed}}{
            % Code for high red
            \stepcounter{highFindings}
        }{
            \ifthenelse{\equal{\currentFindingColor}{mediumOrange}}{
                % Code for medium orange
                \stepcounter{mediumFindings}
            }{
                \ifthenelse{\equal{\currentFindingColor}{lowGreen}}{
                    % Code for low green
                    \stepcounter{lowFindings}
                }{
                    \ifthenelse{\equal{\currentFindingColor}{infoBlue}}{
                        % Code for info blue
                        \stepcounter{informationalFindings}
                    }{
                        % Default case if none match
                    }
                }
            }
        }
    }
}

% importFinding
\newcommand{\importFinding}[1]{%
    \fancyhead[R]{\textcolor{\currentFindingColor}{\Large \findingSection\arabic{finding}}}%
    #1 %
    \closeFinding%
    \incrementFinding%
}

\newcounter{appendixcounter}
\setcounter{appendixcounter}{0}

\newcommand{\newAppendix}[1]{%
    \stepcounter{appendixcounter}
    \subsection{Appendix \Alph{appendixcounter}: #1}
}
%%%

\newcommand{\team}{Finals 10 } % needs trailing space
\newcommand{\teamP}{Finals 10's }
\newcommand{\teamNoSpace}{Finals 10}
\newcommand{\client}{All Ports Tours }
\newcommand{\clientP}{All Ports Tours' }
\newcommand{\clientNoSpace}{All Ports Tours}
\newcommand{\product}{Y }
\newcommand{\productP}{Y's }
\newcommand{\productNoSpace}{Y}


\renewcommand{\texttt}[1]{{\ttfamily #1}}